\section{落地建议：环境先稳定，再谈训练效率}
对初学者来说，Docker 的最大价值不是“更高级”，而是让训练环境可复制。尤其在远程 GPU + 本地 IDE 的工作流里，只要镜像、挂载和工作目录定义清楚，大部分环境问题都能被提前消掉。

\begin{knowledgebox}{推荐的最小工作流}
\begin{itemize}
  \item 先用官方镜像验证依赖是否完整
  \item 再用挂载目录保存代码和实验输出
  \item 最后再接入 VSCode Remote 做日常开发
\end{itemize}
\end{knowledgebox}

\subsection{本章小结}
Docker 不是附属工具，而是 ML 工程的基础设施。先把环境复制问题解决，再去优化训练速度，通常更划算。

\newpage
